% Equations extracted from Cronin (2001)
% New Zealand economic analysis using classical categories

\documentclass{article}
\usepackage{amsmath}

\begin{document}

\section*{Key Equations and Ratios}

% Rate of surplus-value
\subsection*{Rate of Surplus-Value}
\begin{equation}
\frac{s}{v} = \frac{\text{Surplus-Value}}{\text{Variable Capital}}
\end{equation}

% Value composition of capital
\subsection*{Value Composition of Capital}
\begin{equation}
\frac{c}{v} = \frac{\text{Constant Capital}}{\text{Variable Capital}}
\end{equation}

% Accumulation rate
\subsection*{Accumulation Rate}
\begin{equation}
\frac{s}{TV} = \frac{\text{Surplus-Value}}{\text{Total Value}}
\end{equation}

% Total value composition
\subsection*{Total Value}
\begin{equation}
TV = c + v + s
\end{equation}

where:
\begin{itemize}
\item $c$ = Constant Capital (material inputs + depreciation)
\item $v$ = Variable Capital (productive wages)
\item $s$ = Surplus-Value
\item $TV$ = Total Value
\end{itemize}

% Profit-wage ratio (conventional)
\subsection*{Conventional Profit-Wage Ratio}
\begin{equation}
\frac{os}{ec} = \frac{\text{Operating Surplus}}{\text{Employment Compensation}}
\end{equation}

% Rate of profit (implicit from classical framework)
\subsection*{Rate of Profit}
\begin{equation}
r = \frac{s}{c + v}
\end{equation}

This can be rewritten as:
\begin{equation}
r = \frac{s/v}{c/v + 1}
\end{equation}

\end{document}
